\chapter{Blood-Borne Viruses (BBVs)}

\section*{Definition and Overview}
Blood-borne viruses (BBVs) are viruses carried in the blood and body fluids of an infected person that can be passed to others when blood or blood-contaminated fluids enter the bloodstream. The three of greatest concern are the human immunodeficiency virus (HIV), hepatitis B virus (HBV), and hepatitis C virus (HCV). These viruses can cause chronic infection and, over time, serious illness such as AIDS, cirrhosis, and liver cancer. They do not spread through casual contact such as touching, sharing food, or coughing, and each is both preventable and, with modern treatment, manageable. Confidential testing is available through general practice and specialised sexual health and liver services.

\section*{Epidemiology}
Blood-borne viruses affect people in every country. It is estimated that around 40 million people worldwide live with HIV, roughly 250--300 million with chronic hepatitis B, and some 50 million with hepatitis C. In high-income countries, transmission most often occurs through sharing injecting equipment, unprotected sex, and -- in the past -- blood transfusion. Many people with hepatitis B and C do not know they are infected because the conditions can remain silent for years. Vaccination, blood screening, and effective treatment have greatly reduced new infections, and global health strategies aim to eliminate hepatitis B and C and to control HIV.

\section*{Aetiology and Pathophysiology}
Each virus is transmitted when infected blood or body fluid reaches another person's bloodstream:

\begin{itemize}
    \item \textbf{Sharing needles} and injecting equipment for drugs, tattoos, or piercings
    \item \textbf{Sexual contact:} especially for HIV and hepatitis B, through semen, vaginal fluid, and blood
    \item \textbf{From mother to baby:} during pregnancy, birth, or breastfeeding
    \item \textbf{Occupational exposure:} needlestick or splash injuries in healthcare workers
    \item \textbf{Blood transfusion} and organ transplantation, now very rare where blood is screened
    \item \textbf{Hepatitis B and C} infect liver cells directly, causing inflammation that can progress to scarring (cirrhosis) and liver cancer
    \item \textbf{HIV} infects and destroys immune cells, weakening the body's defences against infection and cancer
\end{itemize}

\section*{Clinical Features}
Many people with a BBV have no symptoms at all, especially in the early years. When symptoms occur, they may include:

\begin{itemize}
    \item A flu-like illness shortly after infection (HIV seroconversion or acute hepatitis)
    \item Fatigue and a general feeling of being unwell
    \item Jaundice -- yellowing of the skin and eyes (more typical of hepatitis)
    \item Dark urine and pale stools
    \item Nausea, vomiting, or pain under the ribs
    \item In later HIV: frequent infections, fevers, night sweats, weight loss, and swollen glands
    \item In advanced hepatitis: swelling of the abdomen, easy bruising, and confusion
\end{itemize}

\section*{Differential Diagnosis}
The symptoms of BBV infection overlap with many conditions:

\begin{itemize}
    \item Other viral infections, such as glandular fever, influenza, and other hepatitis viruses
    \item Alcohol-related and non-alcoholic fatty liver disease
    \item Drug-induced liver injury
    \item Autoimmune hepatitis and other liver diseases
    \item Other causes of a flu-like illness or fatigue
    \item Malignancy causing weight loss and gland swelling
\end{itemize}

\section*{When to Seek Medical Care}
Testing is the only way to know if you have a BBV. Seek testing if:

\begin{itemize}
    \item You have shared needles or injecting equipment at any time
    \item You have had sex without a condom, especially with a new or casual partner
    \item You have had a needlestick or blood exposure, including in a healthcare setting
    \item You were born to a mother with HIV or hepatitis
    \item You received a blood transfusion before blood screening was introduced
    \item You have symptoms of hepatitis, such as jaundice or dark urine
\end{itemize}

\textbf{Call 000 (or local emergency services) for:}
\begin{itemize}
    \item Sudden, severe confusion or drowsiness -- a possible sign of advanced liver disease
    \item Vomiting blood or passing black, tarry stools
    \item Difficulty breathing or signs of a severe infection
\end{itemize}

\section*{Diagnosis and Investigations}
Diagnosis is made by blood tests:

\begin{itemize}
    \item \textbf{Screening tests:} blood tests for HIV antibodies or viral particles, and markers of hepatitis B and C
    \item \textbf{Confirmation and monitoring:} further blood tests to confirm infection, measure the viral load, and assess the immune system
    \item \textbf{Liver function tests:} to check how well the liver is working
    \item \textbf{Liver imaging:} ultrasound and newer scans to detect scarring or tumours
    \item \textbf{Liver elastography or biopsy:} to measure the degree of liver scarring in some people with hepatitis
    \item \textbf{Partner and household testing:} so that contacts can also be tested and offered prevention
\end{itemize}

\section*{Management}
All three viruses can now be treated, and in some cases cured:

\begin{itemize}
    \item \textbf{HIV:} antiretroviral therapy is taken lifelong and controls the virus so that people stay well and cannot transmit it
    \item \textbf{Hepatitis C:} modern oral medicines cure the infection in the vast majority of cases, usually over 8--12 weeks
    \item \textbf{Hepatitis B:} long-term antiviral tablets suppress the virus and protect the liver; it is controllable but not usually cured
    \item \textbf{Immunisation:} vaccination protects against hepatitis B, and a vaccine against hepatitis C is in development
    \item \textbf{Supportive care:} healthy eating, limiting alcohol, and looking after the liver
    \item \textbf{Monitoring:} regular blood tests and specialist review for everyone with a BBV
    \item \textbf{Partner notification and counselling:} to prevent further transmission
\end{itemize}

\section*{Complications}
\begin{itemize}
    \item Chronic hepatitis B and C can progress to cirrhosis, liver failure, and liver cancer
    \item Untreated HIV progresses to AIDS, with severe infections and cancers
    \item Liver disease can cause bleeding from the oesophagus, fluid in the abdomen, and confusion
    \item Transmission to partners and babies
    \item Anxiety, stigma, and discrimination
    \item Side effects of treatment, although modern treatments are generally well tolerated
\end{itemize}

\section*{Prognosis and Outcomes}
The outlook has transformed in recent years. Most people with hepatitis C are cured with a short course of tablets, and people with HIV who take treatment can expect a near-normal life span. Chronic hepatitis B is well controlled with daily tablets, and the risk of liver damage falls with effective viral suppression. Untreated, these infections cause serious disease over many years. Early diagnosis and treatment, and avoiding alcohol, give the best outcomes.

\section*{Prevention and Self-Care}
\begin{itemize}
    \item Never share needles, syringes, or other injecting equipment; use needle exchange programs
    \item Use condoms for sex with new or casual partners
    \item Get vaccinated against hepatitis B
    \item Healthcare workers should follow standard precautions against needlestick injuries
    \item Get tested if you may have been exposed, and encourage partners to test
    \item If infected, follow treatment, attend monitoring, and avoid alcohol
    \item Pregnant women are routinely screened so that treatment can protect the baby
    \item Educate yourself and others -- BBVs do not spread through touch, sharing food, or mosquito bites
\end{itemize}

\section*{Sources and Further Reading}
The following sources provide reliable, up-to-date information on blood-borne viruses:

\begin{itemize}
    \item World Health Organization. \textit{HIV, hepatitis B, and hepatitis C fact sheets.} https://www.who.int/
    \item Centers for Disease Control and Prevention. \textit{Viral hepatitis and HIV information.} https://www.cdc.gov/
    \item NHS. \textit{NHS guidance on HIV and viral hepatitis.} https://www.nhs.uk/conditions/
    \item Mayo Clinic. \textit{Information on HIV and viral hepatitis.} https://www.mayoclinic.org/
    \item MedlinePlus. \textit{HIV and hepatitis patient information.} https://medlineplus.gov/
    \item NICE. \textit{HIV and hepatitis testing and management guidance.} https://www.nice.org.uk/
\end{itemize}
